\section{Extracting primes with smooth predecessors}\label{sec:primes}

We prove \Cref{thm:smooth}. Throughout this section $0<\delta<1/4$
is fixed. The task is to detect primes among $N_u:=2u+1$ for a positive
weight on integers $u$ all of whose prime factors are small. The
distribution estimate of \Cref{thm:typeii} will replace tests on a factor
of $N_u$ by simpler tests on rough integers. We specify the order of all
remaining parameter choices at the end of the section.

The extraction separates the composites according to their least prime
factor. A preliminary sieve retains those $N_u$ with no prime factor
below $x^{b_1}$, where $b_1>0$ is fixed and small. For a fixed small
$\kappa>0$ and least prime factors up to $x^{1/2-\kappa}$, Type II
replaces the roughness condition on the complementary factor by a
local density on ordinary integers.
Integrating these densities accounts for the composite part of the
sieve mass. The remaining composites have two prime factors close to
$x^{1/2}$; a separate upper-bound sieve makes their contribution small
with $\kappa$. Its constant must be independent of $\kappa$, so we
will establish that uniformity before making the final choices.

\subsection{The candidate and its mass}

Fix a nonnegative smooth function $\Psi$ supported in $[1,2]$ that is
bounded below by a positive constant on a closed interval of positive
length contained in $(1,2)$. In the marked weight of \Cref{sec:shifts}
take $q=1/2$ and $\ell=1$. Set
\[
 c_1=1,\quad c_2=\frac65,\quad c_3=\frac32,\quad c_4=\frac95,
 \qquad s'=\frac1{r_0(c_1+c_2)},\qquad s_j=s'2^{-j}L,
\]
where the fixed integer $r_0$ is sufficiently large that
\begin{equation}\label{eq:candidate-geometry}
 c_2s'<\delta,\qquad \frac{(r_0-1)c_1}{2}\ge c_2+1.
\end{equation}
Let $j_x$ be the largest integer with $c_1s_{j_x}\ge20T$, and let
\[
 Q_j=\{p\text{ prime}:c_1s_j\le\log p\le c_2s_j\}
 \qquad(0\le j\le j_x).
\]
At each scale for which $[c_3s_j,c_4s_j]$ lies wholly in
$[L^{1/10},L^{1/5}]$ or $[L^{3/10},L^{2/5}]$, take the primes in that
interval of logarithms as a small or big P-group, respectively.
The P-groups and the Q-bands are pairwise disjoint: the four constants
lie in the displayed order, and $c_4<2c_1$. The prime number theorem
and partial summation give fixed upper and positive lower bounds for
each of their reciprocal prime masses. Both numbers of P-groups are
comparable to $T$, and their exponent ranges and endpoint ratios meet
the hypotheses of the preceding sections.

There are $r_0$ ordered slots in each Q-band. Except for one designated
slot in a fixed band $j_{**}$, a slot contains a prime of its band. The
designated slot instead contains any integer $v$ satisfying
\[
 e^{c_1s_{j_{**}}}\le v\le e^{c_2s_{j_{**}}},\qquad \Pminus(v)>W.
\]
We always require $j_{**}\ge j_0$, where $j_0$ is a sufficiently large
fixed lower bound chosen below using only the candidate geometry.
In particular $j_{**}\ge1$. Define
\[
 a_Q(v)=\prod_{j=0}^{j_x}\frac1{r_0!}
       \#\{\text{admissible ordered Q-lists with product }v\},
 \qquad A(u)=a_Q(u_*)W_1(u).
\]
The P-free part $u_*$ and the marked weight $W_1$ are those of
\Cref{sec:shifts}. The rough slot has no P-prime divisor for large $x$,
since all P primes are less than $W$.

The candidate must have substantial ordinary mass near $x$, despite
the restrictions on all its factors. Reciprocal weights make the slot
choices independent. The band geometry leaves room for one prime in
the largest band to adjust their product to size $x$; under this
harmonic sampling, that prime supplies a local probability of order
$1/L$. The next proposition records the resulting mass and the
pointwise bound needed to convert that mass into a count of distinct
primes.

\begin{proposition}[Candidate mass]\label{prop:candidate-mass}
Put
\[
 X_A=\sum_u A(u)\Psi(u/x),\quad
 H_Q=\sum_v\frac{a_Q(v)}v,\quad
 H_P=\sum_{r:r_*=1}\frac{W_1(r)}r.
\]
There are constants $C,C_A$ depending only on the early candidate
choices such that, for sufficiently large $x$ after fixing $j_{**}$,
\begin{align}
 0\le A(u)&\le L^C\tau(u_*),
 & A(u)&\le\exp(C\sqrt L)\quad(u\asymp x),\label{eq:candidate-pointwise}\\
 H_P&\sim\exp\left(q\sum_g V_g\right),
 & X_A&\asymp\frac{x}{L}H_QH_P\gg xL^{-C_A}.
 \label{eq:candidate-mass}
\end{align}
The comparison constants in \eqref{eq:candidate-mass} can be chosen
independently of $j_{**}\ge j_0$. On the support of
$A(u)\Psi(u/x)$ every prime divisor of $u$ lies in $[L^{20},x^\delta]$.
\end{proposition}

\begin{proof}
After fixing the value of the rough slot as a divisor of $v$, the
remaining prime multiplicities determine their band assignments.
The normalized number of ordered lists in each band is at most one.
This gives $a_Q(v)\le\tau(v)$. For $v\ll x$, any possible rough-slot
divisor divides the part of $v$ supported on primes exceeding $W$.
That part has at most $O(L/\log W)=O(\sqrt L)$ prime factors counted
with multiplicity, and hence at most $\exp(O(\sqrt L))$ divisors.
Together with the pointwise bound $W_1(u)\le L^{O(1)}$ this proves
\eqref{eq:candidate-pointwise}. The claimed lower and upper bounds on
prime factors follow from the definitions, \eqref{eq:candidate-geometry},
and $\exp(L^{2/5})<x^{c_2s'}$ for large $x$.

\paragraph{Harmonic masses and the P-part.}
The Q harmonic sum factorizes into the reciprocal masses of all its
slots, divided by $r_0!$ in each band. Each pure-prime slot has mass
between two fixed positive constants. The rough-slot mass is at most
$O(L)$ and is bounded below by a fixed positive constant for
sufficiently large $x$ after fixing $j_{**}$; for the latter assertion
one may restrict that slot to primes of $Q_{j_{**}}$, which then all
exceed $W$. Since there are $O(T)$ slots,
\begin{equation}\label{eq:candidate-harmonic-size}
 L^{-C}\le H_Q\le L^C.
\end{equation}
This exponent and these eventual bounds do not depend on $j_{**}$.

For one P-group write
$F_g(q)=\prod_{p\in\mathcal P_g}(1+q/(p-1))$.
Summing over all powers of its prime divisors shows that its factor
in $H_P$ is
\begin{equation}\label{eq:marked-euler-factor}
 \frac{F_g'(q)}{V_g}
  =F_g(q)\frac1{V_g}\sum_{p\in\mathcal P_g}\frac1{p-1+q}.
\end{equation}
As $\sum_{g,p\in\mathcal P_g}p^{-2}=o(1)$ and $V_g$ is bounded
above and below, multiplying \eqref{eq:marked-euler-factor} proves
the asymptotic for $H_P$. It also proves the useful exact inequality
\begin{equation}\label{eq:mask-euler-comparison}
 H_P\ge\prod_{p\text{ a P prime}}\left(1+\frac q{p-1}\right),
\end{equation}
since $p-1+q\le p$.

The harmonic P-measure assigns negligible mass to $r>x^\sigma$ for
any fixed $\sigma>0$. Indeed set $\xi=L^{-2/5}$ and replace $p^{-1}$
by $p^{-1+\xi}$ in the preceding Euler calculation. Since
$p^\xi\le e$ for every P prime, the tilted factor of each group is
bounded by a fixed constant. Consequently
\begin{equation}\label{eq:candidate-P-tail}
 \sum_{r>x^\sigma,\ r_*=1}\frac{W_1(r)}r
 \le \exp(-\sigma L^{3/5})L^{O(1)}.
\end{equation}

\paragraph{Localizing the product near $x$.}
Normalize the independent P choice and all the Q-slot harmonic
choices by $H_PH_Q$, and denote their product by $U$. Then
\[
 X_A=H_PH_Q\,\E\{U\Psi(U/x)\}.
\]
For the upper comparison we need probability $O(1/L)$ throughout
$x\le U\le2x$; for the lower comparison we need probability
$\gg1/L$ where $U/x$ lies in the positivity interval of $\Psi$.
Both are intervals of bounded length for $\log U$. The P-tail
estimate ensures that the P-part does not move the product out of
reach of the largest Q-band.
Condition on everything except one pure-prime slot in $Q_0$.
The prime number theorem implies that its logarithm falls in any
interval of bounded length with probability $O(1/L)$, uniformly in
the location of that interval. Thus
$\Prob(x\le U\le2x)=O(1/L)$, giving the upper mass bound.

Here is a uniform lower bound. The sum of the midpoints of all the
infinite Q-bands, counting their $r_0$ slots, is exactly
\[
 \sum_{j\ge0}r_0\frac{c_1+c_2}{2}s'2^{-j}L=L.
\]
Let $\Delta=(c_2-c_1)s'L$, the logarithmic width of the free slot.
Choose a fixed prefix of bands so long that the sum of all later
band widths and midpoints is less than $\Delta/24$. Within that
prefix, restrict all slots except the free slot to fixed small
relative neighborhoods of their midpoints, so that their total
deviation is at most $\Delta/24$. This event has probability bounded
below by a positive constant: it involves only a fixed number of
pure-prime slots. Choose $j_0$ larger than the prefix length so that
the exceptional slot is always in the unrestricted tail. Its full
range obeys the same deterministic tail bound. Missing bands after
$j_x$ also obey that bound. Finally restrict the P choice to
$\log r\le\Delta/12$, at negligible loss by
\eqref{eq:candidate-P-tail}.
On this event, a target interval for $\log U$ coming from the
positivity interval of $\Psi$ gives an interval of fixed positive
length for the free slot, lying a distance at least $\Delta/4$
from its range endpoints for large $x$. The prime number theorem
gives conditional probability $\gg1/L$. Multiplying by $U\asymp x$
proves the lower comparison in \eqref{eq:candidate-mass}. This
argument uses no property of the tail distribution beyond its
range, so its constants are uniform in the later choice of
$j_{**}$. The final lower bound follows from
\eqref{eq:candidate-harmonic-size} and the Euler calculation.
\end{proof}

\subsection{Congruence distribution and the preliminary sieve}

For odd squarefree $d$ define
\[
 r(d)=\sum_{d\mid N_u}A(u)\Psi(u/x)-\frac{X_A}{\varphi(d)}.
\]
\begin{proposition}[Type I distribution]\label{prop:typeI}
For every fixed $\vartheta<1/2$ and $D>0$,
\begin{equation}\label{eq:typeI-bound}
 \sum_{\substack{d\le x^\vartheta\\d\text{ odd and squarefree}}}
 |r(d)|\ll xL^{-D}+X_AL^{-18}.
\end{equation}
\end{proposition}

\begin{proof}
The congruence $2u+1\equiv0\pmod d$ specifies a unit class for $u$.
Use character orthogonality on that class. Replacing the principal
coprime mass by $X_A$ has total cost at most $X_AL^{-18}$, because
for each supported $u$,
\[
 \sum_{\substack{d\le x^\vartheta\\(d,u)>1}}
       \frac{\mu^2(d)}{\varphi(d)}
 \le \sum_{p\mid u}\frac1{p-1}
             \sum_{d'\le x}\frac{\mu^2(d')}{\varphi(d')}
 \ll L^{-18}.
\]
Here every $p\mid u$ is at least $L^{20}$, there are at most
$O(L/\log L)$ such primes, and the last sum is $O(L)$ by its Euler
product and Mertens' theorem.

Extract one fixed $Q_0$ slot, as in the proof of
\Cref{prop:candidate-mass}, to write
\[
 A(u)=\sum_{\substack{pr=u\\p\in Q_0}}b(r),
 \qquad 0\le b(r)\le L^{O(1)}\tau(r).
\]
Split $p,r$ into dyads of sizes $P,U$ with $PU\asymp x$.
There is a fixed $a>0$ such that $x^a\ll P,U\ll x^{1-a}$ in
every contributing dyad. A nonprincipal character modulo squarefree
$d$ is induced by a primitive nonprincipal character modulo $f>1$,
where $d=hf$ and $(h,f)=1$. It acts as that primitive character
together with the restrictions $(p,h)=(r,h)=1$, and
$\varphi(d)=\varphi(h)\varphi(f)$.

Fix $h$. On conductors $f\asymp R>L^{C'}$, separate
$\Psi(pr/x)$ by Mellin inversion, which has bounded integrated
absolute cost. For each resulting twist, \Cref{thm:large-sieve},
Cauchy--Schwarz, and \Cref{lem:divisor-moments} give
\begin{align*}
 &\sum_{f\asymp R}\frac1{\varphi(f)}
     \sum_{\chi\bmod f}^{*}
     \left|\sum_{p\asymp P}a_p\chi(p)\right|
     \left|\sum_{r\asymp U}b_r\chi(r)\right|\\
 &\hspace{12mm}\ll
  \frac{L^{O(1)}}R\sqrt{(P+R^2)(U+R^2)PU}\\
 &\hspace{12mm}\ll xL^{O(1)}
       \left(R^{-1}+P^{-1/2}+U^{-1/2}+R/\sqrt x\right).
\end{align*}
The coefficients include coprimality to $h$ and unit-modulus
Mellin twists, so their squared sums are $O(P)$ and
$UL^{O(1)}$. Discarding $(h,f)=1$ and squarefreeness only enlarges
this positive bound. Summing dyads and $1/\varphi(h)$ costs a fixed
power of $L$. Since $R\le x^\vartheta$, a sufficiently large $C'$
gives any prescribed logarithmic saving.

For $1<f\le2L^{C'}$, fix $r$ and apply \Cref{thm:sw} to the
$p$-sum with the original smooth cutoff $\Psi(pr/x)$, using partial
summation. The nonprincipal character sum is $O_A(PL^{-A})$ for
arbitrary fixed $A$. Omitting primes dividing $h$ costs only $O(1)$
terms, since such primes in this range are at least $x^a$ and
$h\le x^\vartheta$. Their total cost after summing $r$, moduli and
dyads is $x^{1-a}L^{O(1)}$. The remaining costs are fixed powers of
$L$, absorbed by choosing $A$ large. This proves
\eqref{eq:typeI-bound}.
\end{proof}

Define
\[
 V(y)=\prod_{p\le y}(1-1/p),\qquad
 V_1(y)=\prod_{3\le p\le y}(1-1/(p-1)),\qquad
 \mathfrak S=2\prod_{p\ge3}\left(1-\frac1{(p-1)^2}\right)>0.
\]
Writing $\gamma_E$ for Euler's constant, Mertens' theorem gives
\begin{equation}\label{eq:sieve-products}
 V(y)\sim\frac{e^{-\gamma_E}}{\log y},\qquad
 V_1(y)\sim\mathfrak S V(y).
\end{equation}
We will choose a small fixed $0<\kappa<1/50$ and then a sufficiently
small fixed $b_1>0$. Set $b_*=b_1/3$ for \Cref{thm:typeii}, and
eventually require
\begin{equation}\label{eq:rough-slot-placement}
 j_{**}\ge j_0,\qquad c_2s'2^{-j_{**}}<b_*/4.
\end{equation}
Apply \Cref{lem:sieve} to odd prime divisors of $N_u$ up to $x^{b_1}$,
with density $g(p)=1/(p-1)$ and the remainders from
\Cref{prop:typeI}. Take
\[
 h=2\left\lfloor\frac1{40b_1}\right\rfloor,
 \qquad \vartheta=\frac12-\frac\kappa2.
\]
For sufficiently small $b_1$, $h$ is large enough for that lemma and
$b_1(4h+2)<\vartheta$. Choose $D>C_A+10$ in
\eqref{eq:typeI-bound}. It follows that
\begin{equation}\label{eq:presieve-mass}
 \sum_{\Pminus(N_u)>x^{b_1}}A(u)\Psi(u/x)
 \ge \frac{\mathfrak S X_A}{L}
   \left\{\frac{e^{-\gamma_E}}{b_1}
                 (1-O(e^{-c/b_1}))+o(1)\right\}.
\end{equation}
The constants in the exponential error are uniform for small
$b_1$ and $0<\kappa<1/50$.

\subsection{Proxies with character and Mellin discrepancy}

We now construct the coefficients that Type II will compare with the
factor tests. The proxy has exactly the same sum as the test on each
short logarithmic cell, but is constant among the $W$-rough integers
in that cell. This count matching handles the principal
characters; prime estimates and the elementary sieve will give the
nonprincipal cancellation independently of Type II. The cell width may
therefore be chosen after Type II specifies its required discrepancy.

At fixed $b_1,\kappa$, partition $(b_1,1/2-\kappa]$ into finitely
many exponent bins $(\gamma,\gamma']$. Their mesh will be chosen
later. On the factor ranges used below consider the tests
\[
 I_\gamma(m)=\ind_{\Pminus(m)>x^\gamma},\qquad
 I_{\mathrm{pr}}(m)=\ind_{m\text{ prime}}.
\]
Every such range lies between two fixed positive powers of $x$;
in particular we can work throughout the ambient interval
$[x^{b_1/4},x^{1-b_1/4}]$. Put $\Delta_L=L^{-S}$, with $S$ a
fixed precision to be chosen last. Partition the logarithmic axis
into intervals of width $\Delta_L$. For a full cell
$C_Y=(Y,e^{\Delta_L}Y]$ meeting this ambient interval, set
\begin{equation}\label{eq:proxy-definition}
 c_I(C_Y)=\frac{\sum_{m\in C_Y}I(m)}
                   {\#\{m\in C_Y:\Pminus(m)>W\}},\qquad
 B_I(m)=c_I(C_Y)\ind_{\Pminus(m)>W}\quad(m\in C_Y).
\end{equation}
The counts defining $c_I$ are full-cell counts even when the factor
range cuts a cell. Range restrictions are imposed afterwards.

\begin{lemma}[Proxy discrepancy]\label{lem:proxy-discrepancy}
The denominators in \eqref{eq:proxy-definition} are positive for
large $x$. Moreover $0\le B_I\le1$. Given any fixed discrepancy
exponent $B$, taking $S>2B+3$ makes
\[
 \alpha_m=(I(m)-B_I(m))\ind_{m\in K}
\]
satisfy \eqref{eq:discrepancy}, uniformly for every interval $K$
inside a dyad in the stated factor ranges, for each of the finitely
many tests above. Thus the coefficient bound required by
\Cref{thm:typeii} is independent of $S$.
\end{lemma}

\begin{proof}
Use \Cref{lem:sieve} on ordinary integers in a cell with $z=W$,
$g(p)=1/p$, and $h=2\lceil T^2\rceil$. Each divisibility remainder
is $O(1)$, and the sieve level is
\[
 D_W=W^{4h+2}=\exp(O(\sqrt L\,T^2))=x^{o(1)}.
\]
Since $Y\gg x^{b_1/4}$, uniformly for these cells,
\begin{equation}\label{eq:proxy-denominator}
 \#\{m\in C_Y:\Pminus(m)>W\}
  =(e^{\Delta_L}-1)YV(W)(1+O(e^{-T^2}))+O(D_W)
  \sim\Delta_L YV(W).
\end{equation}
Both actual tests select subsets of these rough integers for large
$x$: for the prime test all integers in the cell exceed $W$, and
for the other tests $x^\gamma>W$. Hence $0\le c_I\le1$ exactly.

Every prime divisor of a modulus $r\le L^B$ is below $W$, so a
principal character is one throughout the support of $\alpha$.
On a full cell the unweighted sum of $I-B_I$ is zero. For
$f(t)=t^{iv-1}$, $|v|\le L^B$, its variation on that cell is at most
$O((1+|v|)\Delta_L/Y)$. The absolute harmonic mass on a dyad is
$O(1)$ (even the weaker $O(L^{1/2})$ would suffice). Consequently
the full cells contribute $O(L^{B-S})$. Intersecting the test
interval in \eqref{eq:discrepancy} with $K$ leaves at most two
partial cells; their combined harmonic mass is
$O(L^{-S}+x^{-b_1/4})$. These bounds imply the principal-character
case with room to spare when $S>2B+3$.

For a nonprincipal character, first treat the actual roughness
test. Each admitted integer has a unique nondecreasing list of
at most $O(1/b_1)$ prime factors. Fix the first $l-1$ factors and
sum over the last prime. Its order restriction, its lower bound
$x^\gamma$, and the interval restriction on the full product
intersect in one prime interval. Its endpoints lie between
$x^\gamma$ and $x^2$. By \Cref{thm:sw} and partial summation, the
harmonic sum of this prime against a nonprincipal character and
the twist $p^{iv}$ has arbitrarily strong logarithmic saving,
uniformly for a prescribed logarithmic bound on the modulus and
$v$. The reciprocal masses of the remaining factors are bounded
in terms of $b_1$, because
$\sum_{x^\gamma<p\le x^2}p^{-1}=O_{b_1}(1)$.
Summing the bounded number of list lengths preserves the saving.
Repeated primes are included by the nondecreasing enumeration.
Imprimitive nonprincipal characters give the same estimate via
their nonprincipal primitive characters; primes dividing the
modulus are too small to occur. For $I_{\mathrm{pr}}$ this is the
same argument with one prime.

It remains to treat the proxy. For every subinterval $J'$ of a
cell and every unit class $a\pmod r$, apply \Cref{lem:sieve} to
that class, omitting primes dividing $r$ from the sieve. The CRT
gives an $O(1)$ remainder at each squarefree sieve product. Since
all prime divisors of $r$ are below $W$, the main term is
\[
 \frac{|J'|}{r}\prod_{\substack{p\le W\\p\nmid r}}(1-1/p)
 =\frac{|J'|V(W)}{\varphi(r)}.
\]
Thus, also for very short $J'$ as an absolute-error assertion,
\begin{equation}\label{eq:rough-class-count}
 \#\{m\in J':m\equiv a\pmod r,\ \Pminus(m)>W\}
 =\frac{|J'|V(W)}{\varphi(r)}(1+O(e^{-T^2}))+O(D_W).
\end{equation}
Summing against a nonprincipal character cancels the common main
term, and bounds the remaining sum by
$O(|J'|V(W)e^{-T^2}+\varphi(r)D_W)$. Multiply by the cell constant
$c_I\le1$ and use partial summation against $m^{iv-1}$.
The number of cells is at most $O(L^{S+1})$; the moduli and twists
cost fixed powers of $L$. The first error therefore remains
smaller than every fixed negative power of $L$, and the second
is $x^{-b_1/4+o(1)}$ times a fixed power of $L$. This proves the
nonprincipal assertion, including all interval truncations.
\end{proof}

For all later applications, prescribe the saving in
\Cref{thm:typeii} large enough that its errors, after the required
dyadic decompositions, are $o(X_A/L)$. Its coefficient exponent
is fixed before $S$, by \Cref{lem:proxy-discrepancy}; obtain its
required discrepancy exponent and then choose $S$. There are
only finitely many test types, so the same $S$ works for all.

\subsection{The local densities and their integral identity}

The proxy discrepancy allows replacement of a factor test, while its
local density controls the size of the replacement. We need ordinary
rough-integer counts for the cofactor left after a least prime factor
has been selected. For a product of $l$ primes, their logarithms,
divided by $L$, sum to the logarithmic size of that cofactor. This
leads to the following finite sum of integrals over those logarithms.

For $w>\gamma>0$ define
\begin{equation}\label{eq:rough-density-definition}
 D_\gamma(w)=\frac1w+
 \sum_{l\ge2}\frac1{l!}
  \int_{\substack{t_i\ge\gamma\ (1\le i<l)\\
                    \sum_{i<l}t_i\le w-\gamma}}
  \frac1{w-\sum_{i<l}t_i}\prod_{i<l}\frac{dt_i}{t_i}.
\end{equation}
Only finitely many terms are nonzero. On compact subsets of
$w>\gamma>0$ the functions are jointly continuous, including at
thresholds $w=l\gamma$: the integrands have bounded denominators
and the moving boundaries have measure zero.

These are Buchstab's rough-number densities in logarithmic
coordinates, and the minimum-coordinate decomposition below is the
corresponding Buchstab identity \cite{Buchstab1937}.
We derive the short-cell bounds and integral identity needed here directly.

\begin{lemma}[Proxy density bounds]\label{lem:proxy-densities}
For the prime proxy in the interval
$x^{1/2-2\kappa}\le m\le x^{1/2+2\kappa}$,
\begin{equation}\label{eq:prime-proxy-size}
 c_{\mathrm{pr}}(m)\le\frac{C}{LV(W)},
\end{equation}
where $C$ is absolute for $0<\kappa<1/50$. If $mn=N_u$ on the
support of $\Psi(u/x)$ and $\log n/L\in(\gamma,\gamma']$, then
\begin{equation}\label{eq:rough-proxy-size}
 c_\gamma(m)\le
 \frac{\sup_{\gamma\le a\le\gamma'}D_\gamma(1-a)+o(1)}{LV(W)}.
\end{equation}
At fixed $b_1$ one also has
\begin{align}
 \int_{b_1}^{1/2}D_\alpha(1-\alpha)\frac{d\alpha}{\alpha}
       &=D_{b_1}(1)-1,\label{eq:density-integral}\\
 D_{b_1}(1)&\le\frac{e^{-\gamma_E}}{b_1}
                 (1+O(e^{-c/b_1})).\label{eq:density-sieve-upper}
\end{align}
The integrand at the endpoint $\alpha=1/2$ is understood by
its left limit; changing that endpoint value has no effect.
\end{lemma}

\begin{proof}
The prime number theorem with an arbitrarily large logarithmic
error, applied at both endpoints of a cell, gives
\[
 \#\{p\in(Y,e^{\Delta_L}Y]\}
   =(1+o(1))\frac{\Delta_LY}{\log Y}.
\]
For $w=\log Y/L\in[0.46+o(1),0.54+o(1)]$, divide by
\eqref{eq:proxy-denominator} to obtain
\eqref{eq:prime-proxy-size} with an absolute eventual constant.
Increasing fixed $S$ changes the threshold for $x$, not that
constant.

More generally, uniformly on compact sets with $\gamma\ge b_1$,
$w-\gamma$ bounded below positively, and $w$ bounded above, the
number of integers rough above $x^\gamma$ in the cell is at most
\begin{equation}\label{eq:rough-cell-density}
 (D_\gamma(w)+o(1))\frac{\Delta_LY}{L}.
\end{equation}
To see this, repeated prime factors contribute at most
\[
 \sum_{p>x^\gamma}\left\lfloor\frac{2Y}{p^2}\right\rfloor
 \ll Yx^{-b_1}=o(\Delta_LY/L).
\]
Count the squarefree remaining integers by ordered lists of $l$
primes with weight $1/l!$, where $l$ is bounded in terms of
$b_1$. Fix the first $l-1$ primes, write $Q=\prod_{i<l}p_i$ and
$t_i=\log p_i/L$. A last prime can occur only when
$\sum_{i<l}t_i\le w-\gamma+O(\Delta_L/L)$. Dropping its lower
cutoff, its count is at most
\[
 (1+o(1))\frac{\Delta_LY}{LQ(w-\sum_{i<l}t_i)}.
\]
The prime number theorem is uniform here because
$Y/Q\ge x^\gamma/2$, and its logarithmic accuracy is chosen
larger than $S+3$. The reciprocal-prime measures for the remaining
slots converge to $dt_i/t_i$. A finite grid approximation proves
uniform convergence of the resulting integrals: denominators
stay bounded away from zero, while strips around the moving
hyperplane boundary have uniformly vanishing volume. This is
exactly \eqref{eq:rough-density-definition}, proving
\eqref{eq:rough-cell-density}. In the application
$w=1-\log n/L+o(1)$ and $w-\gamma\ge2\kappa+o(1)$.
Continuity, followed by division by
\eqref{eq:proxy-denominator}, proves \eqref{eq:rough-proxy-size}.

For \eqref{eq:density-integral}, write the $l$th term of
$D_{b_1}(1)$, $l\ge2$, on the simplex
\[
 t_1+\cdots+t_l=1,\qquad t_i\ge b_1,
 \qquad \frac1{l!}\frac{dt_1\cdots dt_{l-1}}{t_1\cdots t_l}.
\]
This measure is invariant under every permutation of the $l$
coordinates: exchanging a dependent and an independent coordinate
has absolute Jacobian one. Except on a null set exactly one
coordinate is the minimum, say $\alpha$. Select that coordinate
in $l$ ways. Its measure is $d\alpha/\alpha$, and the remaining
coordinates, with sum $1-\alpha$ and lower bound $\alpha$, give
the $(l-1)$-factor term of $D_\alpha(1-\alpha)$. The coefficient
$l/l!$ is precisely $1/(l-1)!$. Summing proves the identity; the
one-prime term of $D_{b_1}(1)$ equals one.

For the inequality, first obtain a lower density for ordinary
rough integers in $(x,2x]$. Ordered prime lists with weight $1/l!$
assign at most unit mass to every integer, including those with
repetitions. For $l\ge2$ restrict the first $l-1$ logarithms to
$\sum t_i<1-b_1-\eta$, with $\eta>0$ fixed, and sum the last prime
over $(x/Q,2x/Q]$. It is then wholly above the threshold. The
prime number theorem and the same reciprocal-measure convergence
give the corresponding integrals times $x/L$. Include the
one-prime term and let $\eta$ decrease to zero. Null boundaries
give
\[
 \liminf_{x\to\infty}\frac Lx
        \#\{x<m\le2x:\Pminus(m)>x^{b_1}\}\ge D_{b_1}(1).
\]
On the other hand \Cref{lem:sieve}, with $g(p)=1/p$, $O(1)$
remainders and $h=2\lfloor1/(40b_1)\rfloor$, bounds this count by
\[
 xV(x^{b_1})(1+O(e^{-c/b_1}))+O(x^{b_1(4h+2)}).
\]
Its remainder is $o(x/L)$ for small $b_1$, and
\eqref{eq:sieve-products} proves \eqref{eq:density-sieve-upper}.
\end{proof}

The term $1$ in \eqref{eq:density-integral} is the one-prime
contribution to $D_{b_1}(1)$. Selecting a least prime factor accounts
for all the other terms. Thus the integral provides the composite
density to subtract from the preliminary sieve mass, with the
one-prime term left over. These are densities for ordinary rough
integers used in the proxies; their application to $N_u$ uses Type II
followed by Type I and the sieve.

\subsection{Subtracting composites away from balance}

Consider a composite $N_u$ counted in \eqref{eq:presieve-mass},
whose least prime factor is at most $x^{1/2-\kappa}$. Its least
prime factor $p$ lies in one bin $(x^\gamma,x^{\gamma'}]$, and
$m=N_u/p$ satisfies $\Pminus(m)>x^\gamma$. Its contribution is
therefore bounded by
\[
 \sum_{\substack{mp=N_u\\x^\gamma<p\le x^{\gamma'}\text{ prime}}}
       A(u)\Psi(u/x)I_\gamma(m).
\]
All contributing dyads of $m,p$ lie between $x^{b_*}$ and
$x^{1-b_*}$ for large $x$, with $b_*=b_1/3$: the factors bounded
by constants in $N_u\asymp x$ are absorbed by the strict exponent
slack. Apply \Cref{thm:typeii} with
$\alpha=I_\gamma-B_\gamma$ restricted to each actual cofactor
range, and $\beta$ the prime indicator restricted to its bin and
dyad. \Cref{lem:proxy-discrepancy} supplies precisely its
discrepancy hypothesis. Replacing $I_\gamma$ by $B_\gamma$ costs
$o(X_A/L)$ in total.

For a prime $p>W$, the condition that $m=N_u/p$ be rough above
$W$ is equivalent to $\Pminus(N_u)>W$. By
\eqref{eq:rough-proxy-size}, the resulting upper bound reduces
to a constant divided by $LV(W)$ times
\[
 \sum_{x^\gamma<p\le x^{\gamma'}}
       \sum_{\substack{p\mid N_u\\\Pminus(N_u)>W}}
                     A(u)\Psi(u/x).
\]
For each such $p$, sieve the odd primes up to $W$, using the base
mass $X_A/(p-1)$, density $1/(l-1)$, and remainders $r(pd)$.
Take $h=2\lceil T^2\rceil$ in \Cref{lem:sieve}. All products
$pd$ are at most $x^{1/2-\kappa+o(1)}$, which is below the
Type I level $x^{1/2-\kappa/2}$. Moreover different pairs $(p,d)$
give different products: $p>W$ is their unique prime factor above
$W$. Thus \eqref{eq:typeI-bound} bounds the summed remainders
without any divisor multiplicity loss. We obtain
\begin{align*}
 &\sum_{x^\gamma<p\le x^{\gamma'}}
       \sum_{\substack{p\mid N_u\\\Pminus(N_u)>W}}
                  A(u)\Psi(u/x)\\
 &\hspace{12mm}=X_AV_1(W)
          \sum_{x^\gamma<p\le x^{\gamma'}}\frac1{p-1}
                   +o(X_AV(W)).
\end{align*}
Since the reciprocal sum tends to $\log(\gamma'/\gamma)$,
the composite contribution of the bin is at most
\begin{equation}\label{eq:unbalanced-bin}
 \frac{\mathfrak S X_A}{L}
 \left\{\sup_{\gamma\le a\le\gamma'}D_\gamma(1-a)
                    \log(\gamma'/\gamma)+o(1)\right\}.
\end{equation}
At fixed $b_1,\kappa$, joint continuity in
\Cref{lem:proxy-densities} permits a sufficiently fine fixed mesh
such that the sum of these constants is at most
\[
 \int_{b_1}^{1/2-\kappa}D_\alpha(1-\alpha)
                       \frac{d\alpha}{\alpha}+\frac1{10}.
\]
Subtracting \eqref{eq:unbalanced-bin} from
\eqref{eq:presieve-mass} and using
\Cref{eq:density-integral,eq:density-sieve-upper} leaves at least
\begin{equation}\label{eq:prime-balanced-mass}
 \frac{\mathfrak S X_A}{L}
      \left\{1-\frac1{10}
               -O(b_1^{-1}e^{-c/b_1})+o(1)\right\}
 \ge \frac{\mathfrak S X_A}{2L}
\end{equation}
on primes and composites with least prime factor exceeding
$x^{1/2-\kappa}$, provided $b_1$ is sufficiently small.

\subsection{A uniform upper bound for balanced composites}

If a composite $N_u\ll x$ has least prime factor above
$x^{1/2-\kappa}$, it has exactly two prime factors counted with
multiplicity, because $3(1/2-\kappa)>1$. Both factors belong,
for large $x$, to
\[
 \mathcal B_\kappa=[x^{1/2-2\kappa},x^{1/2+2\kappa}].
\]
We may overcount them by ordered pairs $m,n$ of primes in that
range with $mn=N_u$. Replace the first prime indicator by its
proxy using \Cref{thm:typeii}, then replace the second using the
same theorem with the factor roles exchanged. The remaining
coefficient in the second application is a proxy bounded by one.
All required discrepancy and range hypotheses follow from
\Cref{lem:proxy-discrepancy}, as before. By
\eqref{eq:prime-proxy-size}, the resulting upper bound is
$o(X_A/L)$ plus
\begin{equation}\label{eq:balanced-proxy-bound}
 \frac{C}{L^2V(W)^2}
 \sum_{\substack{m,n\in\mathcal B_\kappa,\ mn=2u+1\\
                   \Pminus(m)>W,\ \Pminus(n)>W}}
                     A(u)\Psi(u/x).
\end{equation}
We now bound this positive sum with a constant independent of
$\kappa,b_1,j_{**}$ and $S$.

More precisely, we will show that each dyadic pair with $mn\asymp x$
contributes $O(X_AV(W)^2)$. There are only $O(\kappa L+1)$ such
pairs in the indicated range, so the prefactor in
\eqref{eq:balanced-proxy-bound} will give an
$O((\kappa+L^{-1})X_A/L)$ bound. To prove the dyadic estimate, we
retain a small divisor of $u$ from its Q-slots and P-marks. After
discarding the remaining large Q-prime restrictions, a two-variable
sieve enforces roughness of $m,n$ and the remaining small-prime
restrictions on $u=(mn-1)/2$.

Choose a fixed large index $j'$ and then $\theta$ in the gap
\begin{equation}\label{eq:micro-choice}
 c_2s'2^{-j'}<\theta<c_1s'2^{-j'+1},\qquad
 10\theta\le\frac1{20},\qquad
 \sum_{j\ge j'}r_0c_2s'2^{-j}\le\frac1{25}.
\end{equation}
Such a gap exists because $c_2<2c_1$. These choices depend only
on the candidate geometry. Increase the early lower bound $j_0$
to ensure $j_0\ge j'$, so the rough slot is in a band $j\ge j'$.
Call these the micro bands. Write $H_Q^{\mathrm{mic}}$ for the
Q harmonic mass using just these bands. Omitting the finitely
many fixed bands $j<j'$ gives
\begin{equation}\label{eq:micro-harmonic-comparison}
 H_Q^{\mathrm{mic}}\ll H_Q
\end{equation}
with a constant independent of $j_{**}$.

An assignment consists of the complete ordered Q-lists in the
micro bands, including the rough slot, and one marked prime
from every P-group. Give it weight
\[
 \lambda_{\mathrm{assgn}}
  =\prod_{j=j'}^{j_x}\frac1{r_0!}\prod_g\frac1{V_g},
\]
and let $D_1$ be the product of all its entries, with
multiplicity. The Q contribution has size at most $x^{1/25}$ by
\eqref{eq:micro-choice}; the P marks have product
$\exp(O(TL^{2/5}))=x^{o(1)}$. Thus
\begin{equation}\label{eq:assignment-size}
 D_1\le x^{1/20},\qquad
 \sum_{\mathrm{assgn}}\frac{\lambda_{\mathrm{assgn}}}{D_1}
            =H_Q^{\mathrm{mic}},\qquad
 \sum_{\mathrm{assgn}}\lambda_{\mathrm{assgn}}
            \le x^{1/20}H_Q^{\mathrm{mic}}.
\end{equation}
In the middle identity the harmonic sum of the normalized
P mark in each group is exactly $V_g^{-1}\sum p^{-1}=1$.

The damping in $W_1$ can also be retained in this upper bound.
Once one P-prime per group is marked in $D_1$, every additional
distinct P-prime divisor of $u$ contributes a factor $q$. We represent
these factors by independent random permissions for those primes;
averaging the permissions recovers exactly the damping.

For each assignment, ban every odd non-P prime $l\le x^\theta$
not dividing $D_1$. For each P prime not dividing $D_1$,
independently allow it with probability $q$ and otherwise ban it.
Primes dividing $D_1$ are never banned.
All P primes are below $x^\theta$ and indeed below $W$ for large
$x$. We have the pointwise majorant
\begin{equation}\label{eq:mask-majorant}
 A(u)\le
 \sum_{\mathrm{assgn}}\lambda_{\mathrm{assgn}}
     \ind_{D_1\mid u}\,
     \E_{\mathrm{mask}}\ind_{\{l\nmid u\text{ for every banned }l\}}.
\end{equation}
To verify it, expand $A(u)$ into Q-lists and one P mark per group.
For every genuine list the divisibility $D_1\mid u$ holds, and
every non-P prime divisor of $u$ below $x^\theta$ occurs in
$D_1$: omitted Q-bands contain only pure primes above $x^\theta$.
The micro Q product has no P factors, so $D_1$ contains exactly
one distinct P prime per group. For such a genuine list, the
mask expectation is $q^{\omega(u)-s}$, precisely the damping in
$W_1(u)$. Once the micro entries are fixed, the normalized
number of omitted pure-prime lists with any fixed residual
product is at most one. Explicitly, a band with prime
multiplicities $a_p$ contributes $\prod_p1/a_p!\le1$; disjointness
prevents alternative band assignments. Dropping these residual
restrictions proves \eqref{eq:mask-majorant}.

Fix now a pair of dyads $m\asymp M_1$, $n\asymp M_2$ meeting
the product support $mn\asymp x$, an assignment, and a mask.
Since $D_1$ is odd, $D_1\mid u$ and $mn=2u+1$ imply
$mn\equiv1\pmod{D_1}$. We may drop parity and use this congruence
to bound the number of pairs in the full dyadic box. For an odd
prime $l\le z':=x^\theta$ with $l\nmid D_1$, the bad events are
\[
 mn\equiv0\pmod l\quad(l\le W),\qquad
 mn\equiv1\pmod l\quad(l\text{ banned}).
\]
The two residue sets are disjoint and contain respectively
$2l-1$ and $l-1$ pairs. Put
\[
 \nu(l)=(2l-1)\ind_{l\le W}+(l-1)\ind_{l\text{ banned}},
 \qquad g(l)=\nu(l)/l^2.
\]
These densities satisfy the hypotheses of \Cref{lem:sieve}
uniformly in assignment and mask: $g(l)\le3/l$, and for $l\ge3$
they are bounded away from one. The base mass is the product
of the two real side lengths times $\varphi(D_1)/D_1^2$.

For a squarefree product $d$ of the sieve primes, the CRT gives
exactly $\varphi(D_1)\nu(d)$ residue pairs modulo $D_1d$, including
when $D_1$ has prime-power factors. Each pair of residue classes
has lattice count equal to its area term with error
\[
 O\left(1+\frac{M_1+M_2}{D_1d}\right).
\]
As $\nu(d)\le d\,3^{\omega(d)}$ and $\varphi(D_1)\le D_1$, the
sum of absolute remainders up to $D=(z')^{10}=x^{10\theta}$ is
\begin{equation}\label{eq:pair-lattice-error}
 \ll L^{O(1)}\bigl((M_1+M_2)D+D_1D^2\bigr)
 \le L^{O(1)}\bigl((M_1+M_2)x^{1/20}+D_1x^{1/10}\bigr).
\end{equation}
Here the ordinary $\omega(d)$ occurs only in this elementary
divisor bound; its sums follow, for example, from
$3^{\omega(d)}\le\tau_3(d)$, where $\tau_3(d)$ counts ordered triples
of positive integers with product $d$. Apply the upper-bound version of
\Cref{lem:sieve} with $h=2$.

We record its Euler product explicitly. A banned prime below $W$
has factor $1-3/l+2/l^2$, and a banned prime above $W$ has
factor $1-1/l+1/l^2$. Each is at most
\[
 (1-1/l)^{1+2\ind_{l\le W}}(1+O(l^{-2})).
\]
An allowed P prime has factor $(1-1/l)^2$, so it requires one
compensating factor $(1-1/l)^{-1}$. Omitting the prime 2 changes
only an absolute constant, as does the convergent product of
the $1+O(l^{-2})$ factors. Removed primes dividing $D_1$ require
at most three compensating factors. Therefore
\begin{equation}\label{eq:pair-sieve-product}
 \prod_{\substack{3\le l\le z'\\l\nmid D_1}}(1-g(l))
 \le C V(W)^2V(z')
       \prod_{l\mid D_1}(1-1/l)^{-3}
       \prod_{\substack{p\text{ allowed}\\p\nmid D_1}}(1-1/p)^{-1}.
\end{equation}
Every prime dividing $D_1$ is at least $L^{20}$, so
$D_1\le x^{1/20}$ makes its compensating product uniformly
bounded. The expectation of the last product is
\[
 \prod_{\substack{p\text{ a P prime}\\p\nmid D_1}}
       \left(1-q+\frac q{1-1/p}\right)
 =\prod_{\substack{p\text{ a P prime}\\p\nmid D_1}}
       \left(1+\frac q{p-1}\right)\le H_P
\]
by \eqref{eq:mask-euler-comparison}.

Since $\varphi(D_1)/D_1^2\le1/D_1$, sum the main sieve terms over
assignments using \eqref{eq:assignment-size}, and bound $\Psi$
by its fixed supremum. Their contribution on this dyadic pair is
\[
 \ll xV(W)^2V(x^\theta)H_Q^{\mathrm{mic}}H_P
 \ll X_AV(W)^2,
\]
by \Cref{eq:candidate-mass,eq:micro-harmonic-comparison} and
$V(x^\theta)\ll1/L$. To sum the error
\eqref{eq:pair-lattice-error}, use
$M_i\ll x^{0.54}$ and \eqref{eq:assignment-size}. The result is
\[
 \ll x^{0.64}L^{O(1)}H_Q^{\mathrm{mic}}
       =o(X_AV(W)^2).
\]
The comparison is uniform in the later parameters: use
$H_Q^{\mathrm{mic}}\ll H_Q$, $H_P\ge1$, and
$X_AV(W)^2\asymp xH_QH_P/L^2$.
We have proved the promised dyadic estimate
\begin{equation}\label{eq:balanced-dyadic}
 \sum_{\substack{m\asymp M_1,\ n\asymp M_2,\ mn=2u+1\\
                   \Pminus(m)>W,\ \Pminus(n)>W}}
             A(u)\Psi(u/x)\ll X_AV(W)^2.
\end{equation}

There are $O(\kappa L+1)$ possible dyads for $m$ in
$\mathcal B_\kappa$, and, for each, only $O(1)$ possible dyads
for $n$ because $mn\asymp x$. Combining
\eqref{eq:balanced-proxy-bound} and \eqref{eq:balanced-dyadic}
bounds the entire balanced composite contribution by
\begin{equation}\label{eq:balanced-final}
 C_{\mathrm{bal}}(\kappa+L^{-1})\frac{X_A}{L}+o(X_A/L).
\end{equation}
The constant $C_{\mathrm{bal}}$ depends only on
$\delta,r_0,\Psi,j',\theta$ and the fixed P-group geometry. It
does not depend on $\kappa,b_1,j_{**}$, the later analytic
parameters, or $S$. Indeed the candidate comparison is uniform
in $j_{**}$; the omitted macro bands are fixed before the gap;
the prime proxy bound is uniform on $[0.46,0.54]$; and the
sieve-density and lattice constants just used are absolute.
Thresholds for $x$ may depend on all parameters after they are
fixed.

\subsection{Closing the parameter choices and counting primes}

For clarity, the choices in the preceding argument can be made
in the following order.
\begin{enumerate}
\item Fix $\delta$, $\Psi$, $r_0$ and the Q/P geometry. Choose the
      finite prefix for the candidate lower bound, then $j'$ and
      $\theta$ in \eqref{eq:micro-choice}, and enlarge $j_0$ to
      cover both requirements. All these are early choices.
\item Fix $\kappa\in(0,1/50)$ so small that
      $C_{\mathrm{bal}}\kappa<\mathfrak S/4$. This is possible
      because the constant in \eqref{eq:balanced-final} is
      independent of the gap.
\item Choose $b_1$ sufficiently small for
      \eqref{eq:presieve-mass}, \eqref{eq:density-sieve-upper},
      and \eqref{eq:prime-balanced-mass}. Set $b_*=b_1/3$.
      Choose $j_{**}$ to meet \eqref{eq:rough-slot-placement},
      and then choose the finite exponent mesh for
      \eqref{eq:unbalanced-bin}.
\item Prescribe the saving in \Cref{thm:typeii} so that every
      replacement error, including dyadic sums, is $o(X_A/L)$.
      The coefficient exponents are already bounded independently
      of $S$. The proof of that theorem fixes its Cauchy splitting
      precision, shift-saving target, graph and analytic
      parameters, and finally a required discrepancy exponent
      in \eqref{eq:discrepancy}.
\item Choose $S$ sufficiently large in
      \Cref{lem:proxy-discrepancy} for that exponent, fix any
      prime-number-theorem accuracies needed for this $S$, and
      let $x$ tend to infinity.
\end{enumerate}
In particular the precision of the proxy cells causes no
feedback into the Type II coefficient bound.

Subtract \eqref{eq:balanced-final} from
\eqref{eq:prime-balanced-mass}. For all sufficiently large $x$
there remains a positive constant multiple of $X_A/L$ on prime
values of $N_u$:
\[
 \sum_{2u+1\text{ prime}}A(u)\Psi(u/x)\gg X_A/L.
\]
By \Cref{prop:candidate-mass} each summand is at most
$\exp(O(\sqrt L))$ and $X_A\gg xL^{-C_A}$. The map
$u\mapsto2u+1$ is injective, so the number of distinct primes is
at least
\[
 xL^{-C_A-1}\exp(-O(\sqrt L))=x^{1-o(1)}.
\]
Their range is $2x<p\le4x+1\le5x$. Finally
$\Pplus(p-1)=\Pplus(2u)=\max(2,\Pplus(u))\le x^\delta$ for
large $x$; the strict inequality $c_2s'<\delta$ in the candidate
geometry leaves room for every subpower P prime, and the factor
2 is eventually below $x^\delta$. This proves
\Cref{thm:smooth}.
